% Unit 048 terjemahan Bahasa Indonesia dari chapter4.tex baris 9--60.
% Sumber: Methods of Algebra, Volume 2, commit
% 9a5803ff77dd3257484cb177f851a73770a59dd3 (CC BY 4.0).
% stable-unit-id: o014.aljabr2.chapter4.overview
% Cakupan: ikhtisar Bab 4 sebelum Bagian 4.1.

% segment-id: o014.li2.chapter4.overview.h001
\chapter{Kategori Bertriangulasi dan Kategori Turunan\texorpdfstring{\hspace{0.1em}}{}}%
\label{sec:triangulated-derived-cat}
\hypertarget{o014.aljabr2.chapter4.overview}{ }

% segment-id: o014.li2.chapter4.overview.p001
Kategori turunan dirintis oleh A.\ Grothendieck dan J.-L.\ Verdier pada
dasawarsa 1960-an untuk mempelajari teorema dualitas dalam geometri aljabar.
Perkembangan selanjutnya menunjukkan bahwa kategori turunan menyediakan
bahasa yang praktis dan ampuh, baik dalam geometri, topologi, maupun teori
representasi. Informasi yang terkandung dalam kategori turunan
$\cate{D}(\mathcal{A})$ dari kategori abelian $\mathcal{A}$ berada di antara
kompleks dan kohomologinya. Informasi tersebut diungkapkan secara tepat melalui
kategori bertriangulasi yang diperkenalkan dalam
\S\S~\sourcecrossref{sec:triangular-def}{4.1}--%
\sourcecrossref{sec:triangular-basic}{4.2}: $\cate{D}(\mathcal{A})$ dilengkapi
autoekuivalensi pergeseran $X\mapsto X[1]$ serta suatu kelas barisan morfisme
yang disebut segitiga terbedakan (juga disebut segitiga eksak),
% segment-id: o014.li2.chapter4.overview.q001
\[
  X \xrightarrow{f} Y \xrightarrow{g} Z \xrightarrow{h} X[1],
  \quad \text{atau disingkat}\quad X\to Y\to Z\xrightarrow{+1} .
\]
Segitiga-segitiga tersebut memenuhi aksioma (TR0)--(TR5). Sebagai contoh,
suatu segitiga terbedakan dapat dirotasi menjadi
% segment-id: o014.li2.chapter4.overview.q002
\[
  Y \xrightarrow{g} Z \xrightarrow{h} X[1] \xrightarrow{-f[1]} Y[1]
  \quad\text{atau}\quad
  Z[-1] \xrightarrow{-h[-1]} X \xrightarrow{f} Y \xrightarrow{g} Z .
\]

% segment-id: o014.li2.chapter4.overview.p002
Segitiga terbedakan menggantikan barisan eksak pendek dalam kategori kompleks
$\cate{C}(\mathcal{A})$, meskipun analogi yang lebih tepat ialah dengan
barisan kofiber dalam teori homotopi. Funktor kohomologis dari kategori
bertriangulasi ke kategori abelian menyediakan mekanisme untuk menghasilkan
barisan eksak panjang dari segitiga terbedakan.

% segment-id: o014.li2.chapter4.overview.p003
Peralihan dari $\cate{C}(\mathcal{A})$ ke $\cate{D}(\mathcal{A})$ merupakan
pokok bahasan \S~\sourcecrossref{sec:derived-cat}{4.4} dan berlangsung dalam
dua langkah.
% segment-id: o014.li2.chapter4.overview.l001
\begin{enumerate}
	% segment-id: o014.li2.chapter4.overview.i001
	\item Langkah pertama beralih dari $\cate{C}(\mathcal{A})$ ke
	$\cate{K}(\mathcal{A})$. Kategori ini dilengkapi funktor pergeseran
	$X\mapsto X[1]$ dan struktur kategori bertriangulasi. Untuk setiap
	morfisme $f:X\to Y$, kerucut pemetaannya beserta morfisme kanonik
	% segment-id: o014.li2.chapter4.overview.q003
	\[
	  X \xrightarrow{f} Y \xrightarrow{\alpha(f)} \Cone(f)
	  \xrightarrow{\beta(f)} X[1]
	\]
	membentuk segitiga terbedakan yang bersesuaian hingga isomorfisme. Perangkat
	untuk langkah ini telah disiapkan pada bagian pertama \CHref{sec:cplx}.
	% segment-id: o014.li2.chapter4.overview.i002
	\item Langkah kedua secara formal menambahkan invers bagi semua
	kuasi-isomorfisme, atau secara ekuivalen menolkan semua kompleks asiklik,
	melalui lokalisasi. Hasilnya ialah $\cate{D}(\mathcal{A})$, yang mempunyai
	kumpulan objek yang sama dengan $\cate{C}(\mathcal{A})$ dan mewarisi
	struktur kategori bertriangulasi dari $\cate{K}(\mathcal{A})$. Langkah ini
	bertumpu pada teori umum lokalisasi Verdier dalam
	\S~\sourcecrossref{sec:triangulated-localization}{4.3}. Di sinilah aksioma
	oktahedral (TR5), aksioma paling rumit dalam definisi kategori
	bertriangulasi, menjalankan perannya.
\end{enumerate}

% segment-id: o014.li2.chapter4.overview.p004
Mengambil kohomologi kompleks menghasilkan funktor kohomologis pada tingkat
kategori turunan, $\Hm^0:\cate{D}(\mathcal{A})\to\mathcal{A}$, dengan
$\Hm^n(X):=\Hm^0(X[n])$. Jika proses yang sama dimulai dari kategori kompleks
terbatas bawah (atau terbatas atas, atau terbatas), diperoleh kategori
bertriangulasi $\cate{D}^+(\mathcal{A})$ (atau
$\cate{D}^-(\mathcal{A})$, atau $\cate{D}^{\bdd}(\mathcal{A})$). Kategori
tersebut juga dapat dipandang sebagai subkategori bertriangulasi penuh dari
$\cate{D}(\mathcal{A})$, yang masing-masing dicirikan oleh syarat
$n\ll0$ (atau $n\gg0$, atau $|n|\gg0$) $\implies\Hm^n(X)=0$.

% segment-id: o014.li2.chapter4.overview.p005
Secara lebih umum, syarat lenyapnya kohomologi dapat dipakai untuk
mendefinisikan subkategori penuh $\cate{D}^{\geq0}(\mathcal{A})$ dan
$\cate{D}^{\leq0}(\mathcal{A})$ dari $\cate{D}(\mathcal{A})$, dan funktor
pemenggalan kompleks tetap dapat didefinisikan pada kategori turunan.
Funktor kanonik $\mathcal{A}\to\cate{D}(\mathcal{A})$ bersifat setia penuh dan
mengidentifikasi $\mathcal{A}$ dengan
$\cate{D}^{\leq0}(\mathcal{A})\cap\cate{D}^{\geq0}(\mathcal{A})$. Jika
$\mathcal{A}$ mempunyai cukup banyak objek injektif atau projektif, maka
$\Hom_{\cate{D}(\mathcal{A})}(X,Y[n])\simeq
\Ext^n_{\mathcal{A}}(X,Y)$, dengan $X$ dan $Y$ objek $\mathcal{A}$ yang
dipandang sebagai kompleks terkonsentrasi di derajat nol.%
\footnote{Catatan penerjemah (O014-C071): sumber menyebut $X$ dan $Y$ sebagai
kompleks pada $\mathcal{A}$. Namun, penyematan setia penuh
$\mathcal{A}\to\cate{D}(\mathcal{A})$ pada kalimat sebelumnya dan definisi
$\Ext^n_{\mathcal{A}}$ dalam \S~\ref{sec:Ext-Tor} memberikan rumus ini bagi
objek $\mathcal{A}$ yang dipandang sebagai kompleks derajat nol; target
memulihkan cakupan tersebut.}
Definisi $\Ext^n_{\mathcal{A}}$ diberikan dalam \S~\ref{sec:Ext-Tor}, dan
semua hal di atas dibahas dalam \S~\sourcecrossref{sec:Hom-Ext}{4.5}.

% segment-id: o014.li2.chapter4.overview.p006
Untuk funktor eksak kiri $F:\mathcal{A}\to\mathcal{A}'$, funktor turunan
kanan klasiknya terwujud pada tingkat kategori turunan terbatas bawah sebagai
$\mathrm{R}F:\cate{D}^+(\mathcal{A})\to\cate{D}^+(\mathcal{A}')$. Funktor
ini dapat dicirikan sebagai ekstensi Kan kiri; data terkait dituliskan dengan
diagram $2$-sel berikut.
% segment-id: o014.li2.chapter4.overview.g001
\[\begin{tikzcd}
	\cate{K}^+(\mathcal{A}) \arrow[r, "{\cate{K}F}"] \arrow[d, "Q"'] &
	\cate{K}^+(\mathcal{A}') \arrow[d, "{Q'}"] \arrow[Rightarrow, ld] \\
	\cate{D}^+(\mathcal{A}) \arrow[r, "{\mathrm{R}F}"'] &
	\cate{D}^+(\mathcal{A}')
\end{tikzcd}\]

Di sini $Q$ dan $Q'$ menyatakan funktor lokalisasi, sementara
$\mathrm{R}F$ juga harus dilengkapi struktur funktor bertriangulasi. Funktor
klasik $\mathrm{R}^nF$ diperoleh dengan membatasi $\mathrm{R}F$ pada
$\mathcal{A}$ lalu mengambil $\Hm^n$. Jika $F$ justru funktor eksak kanan,
pengangkatan funktor turunan kiri klasik $\mathrm{L}_nF$ dituliskan sebagai
berikut.%
\footnote{Catatan penerjemah (O014-C070): pada diagram berikut sumber memberi
label $Q$ pada panah vertikal kanan, sedangkan objek sumber dan targetnya
masing-masing ialah $\cate{K}^-(\mathcal{A}')$ dan
$\cate{D}^-(\mathcal{A}')$. Sesuai diagram sebelumnya dan penamaan dalam
prosa, target memulihkan label $Q'$.}
% segment-id: o014.li2.chapter4.overview.g002
\[\begin{tikzcd}
	\cate{K}^-(\mathcal{A}) \arrow[r, "{\cate{K}F}"] \arrow[d, "Q"'] &
	\cate{K}^-(\mathcal{A}') \arrow[d, "{Q'}"] \\
	\cate{D}^-(\mathcal{A}) \arrow[r, "{\mathrm{L}F}"']
	\arrow[ru, Rightarrow] & \cate{D}^-(\mathcal{A}').
\end{tikzcd}\]

% segment-id: o014.li2.chapter4.overview.p007
Keberadaan funktor turunan dan struktur terperinci kategori turunan, sebagaimana
dalam teori klasik, bergantung pada resolusi kompleks. Resolusi menyediakan
perangkat yang diperlukan untuk bekerja secara konkret dengan kategori turunan.
Rinciannya terdapat dalam
\S\S~\sourcecrossref{sec:triangulated-functor-localization}{4.6}--%
\sourcecrossref{sec:bounded-derived-functor}{4.8}. Versi tak terbatas
$\cate{D}(\mathcal{A})\to\cate{D}(\mathcal{A}')$ mempunyai rumusan serupa
dan akan dibahas dalam \S~\sourcecrossref{sec:unbounded-derived}{4.11}.

% segment-id: o014.li2.chapter4.overview.p008
Dengan mempertimbangkan amplitudo funktor turunan, yakni ``pergeseran'' yang
dihasilkannya pada tingkat kohomologi, kita dapat membahas dimensi funktor
turunan (Definisi--Proposisi~\sourcecrossref{def:dim-functor}{pada \S~4.7}).
Di bawah syarat yang sesuai, pembentukan funktor turunan juga kompatibel dengan
komposisi funktor
(Teorema~\sourcecrossref{prop:localization-injective}{pada \S~4.8}). Hal ini
memberikan penjelasan seragam bagi sejumlah barisan spektral dalam teori
klasik, yang disebut barisan spektral Grothendieck dan dijabarkan dalam
Teorema~\sourcecrossref{prop:Grothendieck-ss}{pada \S~4.8}. Teori umum
barisan spektral ditangguhkan sampai Bab~\sourcecrossref{sec:ss}{5}.

% segment-id: o014.li2.chapter4.overview.p009
Seluruh konstruksi ini dapat dirumuskan untuk kategori bertriangulasi umum dan
lokalisasi Verdier-nya. Perlu dicatat bahwa dalam penerapan sering dijumpai
kategori bertriangulasi yang tidak berasal langsung dari kategori abelian,
atau funktor kohomologis yang bukan $\Hm^n$. Dalam situasi umum itu, peran
funktor pemenggalan serta $\cate{D}^{\leq0}$ dan $\cate{D}^{\geq0}$ digantikan
oleh struktur-$t$ pada kategori bertriangulasi, tetapi buku ini tidak membahas
topik tersebut.

% segment-id: o014.li2.chapter4.overview.p010
Funktor turunan yang dibahas secara konkret dalam bab ini dibatasi pada contoh
yang muncul dalam aljabar. Contoh pertama ialah funktor $\RHom$ dalam
\S~\sourcecrossref{sec:RHom}{4.9}. Funktor ini merupakan versi kategori
turunan dari funktor $\Ext$ dan melibatkan teori bifunktor turunan. Sejumlah
keadjungan juga akan ditingkatkan ke kategori turunan dengan prinsip mengganti
$\Hom$ oleh $\RHom$. Berikutnya,
\S~\sourcecrossref{sec:Rlim}{4.10} membahas hasil kali dan koproduk dalam
$\cate{D}(\mathcal{A})$, kemudian memperkenalkan pengangkatan $\lim^n$
yang ditulis $\mathrm{R}\lim$. Seperti $\lim^n$, funktor ini mempunyai
deskripsi ringkas melalui limit homotopi
(Definisi~\sourcecrossref{def:holim}{pada \S~4.10}), meskipun pada tingkat
kategori turunan deskripsi tersebut hanya dapat diberikan secara kasar.

% segment-id: o014.li2.chapter4.overview.p011
Selanjutnya, \S~\sourcecrossref{sec:otimesL}{4.12} memperkenalkan versi
kategori turunan dari funktor $\Tor$, yaitu hasil kali tensor turunan
$\otimesL$, suatu bifunktor bertriangulasi. Bagian itu juga membuktikan banyak
fakta, termasuk kendala asosiativitas
(Proposisi~\sourcecrossref{prop:otimesL-associativity-constraint}{pada
\S~4.12}) dan keadjungan versi $\RHom$
(Teorema~\sourcecrossref{prop:RHom-otimesL-adjunction}{pada \S~4.12}). Seperti
dalam teori klasik, $\RHom$ dan $\otimesL$ merupakan dua contoh dasar funktor
turunan. Agar deskripsinya lengkap dan bersih, bagian tersebut memakai
kategori turunan tak terbatas beserta funktor turunannya. Karena itu diperlukan
teori dalam \S~\sourcecrossref{sec:unbounded-derived}{4.11}, khususnya konsep
resolusi K-projektif dan K-datar; dalam konteks teori modul yang dibahas di
sini, keberadaan resolusi tersebut tidak menjadi masalah.

% segment-id: o014.li2.chapter4.overview.p012
Terakhir, perlu dicatat beberapa kekurangan kategori turunan
$\cate{D}(\mathcal{A})$.
% segment-id: o014.li2.chapter4.overview.l002
\begin{enumerate}
	% segment-id: o014.li2.chapter4.overview.i003
	\item Dalam aksioma kategori bertriangulasi, (TR2) menyatakan bahwa setiap
	morfisme dapat ditempatkan dalam suatu segitiga terbedakan. Segitiga
	semacam itu unik hingga isomorfisme, tetapi bukan ``hingga isomorfisme
	unik'' sebagaimana lazim dalam aljabar.
	% segment-id: o014.li2.chapter4.overview.i004
	\item Kategori bertriangulasi biasanya tidak mempunyai kernel maupun
	kokernel, sehingga operasi $\varinjlim$ dan $\varprojlim$ sangat sulit
	dijalankan. Kerucut pemetaan menyediakan pengganti dalam arti homotopi,
	tetapi, seperti telah disebutkan, keunikannya dalam kategori turunan
	bermasalah.
	% segment-id: o014.li2.chapter4.overview.i005
	\item Meskipun funktor turunan pada tingkat kategori turunan mempunyai
	definisi yang jelas, tidak seperti funktor-$\delta$ kohomologis universal,
	funktor tersebut sulit dicirikan secara sederhana melalui sifat
	keterhapusan;
	lihat Proposisi~\ref{prop:erasable-univ-delta}.
\end{enumerate}

% segment-id: o014.li2.chapter4.overview.p013
Akar persoalannya ialah bahwa dalam membangun $\cate{D}(\mathcal{A})$, relasi
homotopi dihapus secara terlalu kasar. Kekurangan ini perlu diatasi dengan
konstruksi yang lebih tinggi daripada kategori turunan dan mungkin juga lebih
alami; kategori turunan semestinya hanyalah bayangan dari konstruksi tersebut.

% segment-id: o014.li2.chapter4.overview.n001
\begin{petunjukbacaan}
	Sebagian besar isi bab ini merupakan konsep inti dalam teori kategori
	bertriangulasi dan kategori turunan. Materi sampingan mencakup pembahasan
	ekstensi-$n$ dalam \S~\sourcecrossref{sec:Hom-Ext}{4.5}, terutama yang
	berpusat pada Teorema Yoneda
	(Teorema~\sourcecrossref{prop:Yoneda-Ext}{pada \S~4.5}), serta pembahasan
	funktor turunan tak terbatas dalam
	\S~\sourcecrossref{sec:unbounded-derived}{4.11}, yang bertumpu pada teori
	\S~\ref{sec:K-injectives}. Kedua topik ini merupakan pengetahuan umum bagi
	para spesialis, tetapi pembaca hendaknya menilai sendiri kebutuhannya.
	Jika pembaca hanya membahas funktor turunan antara kategori turunan
	terbatas bawah atau terbatas atas, semua pernyataan dalam
	\S~\sourcecrossref{sec:otimesL}{4.12} mengenai resolusi K-projektif (atau
	K-datar) dapat diganti dengan resolusi projektif (atau datar) dari kompleks
	terbatas atas tanpa mengganggu alur bacaan. Keabsahan penggantian ini
	dijamin oleh Proposisi~\sourcecrossref{prop:flat-K-flat}{pada \S~4.11}.
\end{petunjukbacaan}
